%【紙面設定】
%b4,横書き,2段
\documentclass[paper=b4j,landscape,twocolumn,fleqn]{jlreq}
%横書き用の調整
\special{papersize=\the\paperwidth,\the\paperheight}
%間に線をいれる
\setlength{\columnseprule}{0.4pt}
%ページ番号消す
\pagestyle{empty}
%ページ余白の調整
\usepackage[margin=8truemm]{geometry}
%画像挿入用のパッケージ
\usepackage[dvipdfmx]{graphicx}
%数学用のパッケージ
\usepackage{amsmath}
%見出しによる行取りの数を減らす
\ModifyHeading{section}{
	lines=1,
	before_space=6pt,
	after_space=0pt
}
% 段落・数式まわりの余白を詰める
\setlength{\parskip}{0pt}
\setlength{\jot}{3pt}

%【本文】
\begin{document}
\normalsize
%1段落目
\section{成分表示と大きさ}
次のベクトルについて、成分表示・大きさ・単位ベクトルを求めよ。
\begin{align*}
　&1.　\vec{a}=\begin{pmatrix}3 & 4\end{pmatrix}\text{ の大きさ }|\vec{a}|\text{ を求めよ。}\\[0.45em]
　&2.　\vec{b}=\begin{pmatrix}-5 & 12\end{pmatrix}\text{ の大きさ }|\vec{b}|\text{ を求めよ。}\\[0.45em]
　&3.　\vec{c}=\begin{pmatrix}6 & 8\end{pmatrix}\text{ の単位ベクトル }\hat{c}\text{ を求めよ。}\\[0.45em]
　&4.　\vec{d}=\begin{pmatrix}-4 & 3\end{pmatrix}\text{ の単位ベクトル }\hat{d}\text{ を求めよ。}\\[0.45em]
　&5.　長さが5で、方向が \begin{pmatrix}3 & 4\end{pmatrix} と同じベクトルを求めよ。\\[0.45em]
　&6.　\vec{v}=\begin{pmatrix}x & y\end{pmatrix}\text{ の大きさが10のとき、成り立つ式を立てよ。}\\[0.45em]
\end{align*}

\section{点と点を結ぶベクトル}
点A, Bから、\(\overrightarrow{AB}\) と距離 \(\left|\overrightarrow{AB}\right|\) を求めよ。
\begin{align*}
　&1.　A(1,2),\ B(5,7)\\[0.45em]
　&2.　A(-3,4),\ B(2,-2)\\[0.45em]
　&3.　A(0,0),\ B(-6,8)\\[0.45em]
　&4.　A(4,-1),\ B(-2,5)\\[0.45em]
　&5.　A(2,3),\ B(2,-5)\\[0.45em]
　&6.　A(a,b),\ B(c,d)\text{ のとき、}\overrightarrow{AB}\text{ と }|\overrightarrow{AB}|\text{ を一般式で表せ。}\\[0.45em]
\end{align*}

\section{ベクトルの足し算・引き算・スカラー倍}
各問で計算し、結果を成分で答えよ。
\begin{align*}
　&1.　\begin{pmatrix}2 & -1\end{pmatrix}+\begin{pmatrix}5 & 3\end{pmatrix}\\[0.45em]
　&2.　\begin{pmatrix}-4 & 7\end{pmatrix}-\begin{pmatrix}1 & -2\end{pmatrix}\\[0.45em]
　&3.　2\begin{pmatrix}3 & -5\end{pmatrix}\\[0.45em]
　&4.　-3\begin{pmatrix}-2 & 4\end{pmatrix}\\[0.45em]
　&5.　\begin{pmatrix}1 & 2\end{pmatrix}+\begin{pmatrix}-3 & 5\end{pmatrix}-\begin{pmatrix}4 & -1\end{pmatrix}\\[0.45em]
　&6.　2\vec{a}-3\vec{b}\text{ を、}\vec{a}=\begin{pmatrix}1 & -4\end{pmatrix},\ \vec{b}=\begin{pmatrix}-2 & 3\end{pmatrix}\text{ のとき求めよ。}\\[0.45em]
\end{align*}

\newpage

\section{内積と角度}
次のベクトルについて、内積となす角を求めよ。
\begin{align*}
　&1.　\vec{a}=\begin{pmatrix}1 & 2\end{pmatrix},\ \vec{b}=\begin{pmatrix}3 & 4\end{pmatrix}\text{ の内積 }\vec{a}\cdot\vec{b}\\[0.45em]
　&2.　\vec{a}=\begin{pmatrix}2 & -1\end{pmatrix},\ \vec{b}=\begin{pmatrix}1 & 2\end{pmatrix}\text{ の内積 }\vec{a}\cdot\vec{b}\\[0.45em]
　&3.　\vec{a}=\begin{pmatrix}3 & 0\end{pmatrix},\ \vec{b}=\begin{pmatrix}0 & 5\end{pmatrix}\text{ が垂直か判定せよ。}\\[0.45em]
　&4.　\vec{a}=\begin{pmatrix}2 & 1\end{pmatrix},\ \vec{b}=\begin{pmatrix}1 & 2\end{pmatrix}\text{ のなす角 }\theta\text{ を求めよ。}\\[0.45em]
　&5.　\vec{a}=\begin{pmatrix}4 & 0\end{pmatrix},\ \vec{b}=\begin{pmatrix}2 & 2\sqrt{3}\end{pmatrix}\text{ のなす角 }\theta\text{ を求めよ。}\\[0.45em]
　&6.　\vec{a}=\begin{pmatrix}4 & -1\end{pmatrix},\ \vec{b}=\begin{pmatrix}1 & 4\end{pmatrix}\text{ について、}\vec{a}\cdot\vec{b}\text{ を計算し、2本が垂直か判定せよ。}\\[0.45em]
\end{align*}

\section{2次元の外積と面積}
次のベクトルについて、外積(2次元)を計算せよ。
\begin{align*}
　&1.　\vec{a}=\begin{pmatrix}1 & 2\end{pmatrix},\ \vec{b}=\begin{pmatrix}3 & 1\end{pmatrix}\\[0.45em]
　&2.　\vec{a}=\begin{pmatrix}2 & 5\end{pmatrix},\ \vec{b}=\begin{pmatrix}-1 & 4\end{pmatrix}\\[0.45em]
　&3.　\vec{a}=\begin{pmatrix}4 & 0\end{pmatrix},\ \vec{b}=\begin{pmatrix}0 & 3\end{pmatrix}\\[0.45em]
　&4.　問3の2本が作る平行四辺形の面積を求めよ。\\[0.45em]
　&5.　\vec{a}\times\vec{b}\text{ と }\vec{b}\times\vec{a}\text{ の関係を式で書け。}\\[0.45em]
　&6.　\vec{a}=\begin{pmatrix}2 & 1\end{pmatrix},\ \vec{b}=\begin{pmatrix}1 & 3\end{pmatrix}\text{ と }\vec{c}=\begin{pmatrix}3 & -1\end{pmatrix}\text{ について、}\vec{a}\times\vec{b}\text{ と }\vec{a}\times\vec{c}\text{ を計算し、符号を答えよ。}\\[0.45em]
\end{align*}

\section{射影とゲーム計算}
ベクトル\(\vec{a}\)から\(\vec{b}\)方向への射影を\(\vec{A}\)とする。
\begin{align*}
　&1.　\vec{a}=\begin{pmatrix}3 & 2\end{pmatrix},\ \hat{b}=\begin{pmatrix}1 & 0\end{pmatrix}\text{ のとき、}\vec{A}\text{ を求めよ。}\\[0.45em]
　&2.　\vec{a}=\begin{pmatrix}4 & 1\end{pmatrix},\ \hat{b}=\begin{pmatrix}0 & 1\end{pmatrix}\text{ のとき、}\vec{A}\text{ を求めよ。}\\[0.45em]
　&3.　\vec{a}=\begin{pmatrix}5 & 0\end{pmatrix},\ \hat{b}=\begin{pmatrix}\frac{3}{5} & \frac{4}{5}\end{pmatrix}\text{ のとき、}\vec{A}\text{ を求めよ。}\\[0.45em]
　&4.　点A(2,2),\ 点B(5,3.6),\ 半径r_a=1.2,\ r_b=1.8\text{ とする。衝突判定せよ。}\\[0.45em]
　&5.　壁方向 \hat{w}=\begin{pmatrix}\frac{1}{\sqrt{2}} & \frac{1}{\sqrt{2}}\end{pmatrix},\ 入力\vec{v}=\begin{pmatrix}4 & 0\end{pmatrix}\text{ のとき、壁沿い移動 }\vec{v}_{\parallel}\text{ を求めよ。}\\[0.45em]
　&6.　壁方向 \hat{w}=\begin{pmatrix}0 & 1\end{pmatrix},\ 入力\vec{v}=\begin{pmatrix}3 & -2\end{pmatrix}\text{ のとき、壁沿い移動 }\vec{v}_{\parallel}\text{ を求めよ。}\\[0.45em]
\end{align*}

\end{document}
